% Chương 2 — Nền tảng xử lý ngôn ngữ tiếng Việt. Nạp bằng % Chương 2 — Nền tảng xử lý ngôn ngữ tiếng Việt. Nạp bằng % Chương 2 — Nền tảng xử lý ngôn ngữ tiếng Việt. Nạp bằng % Chương 2 — Nền tảng xử lý ngôn ngữ tiếng Việt. Nạp bằng \input{chapter2_nlp}.
% Chỉ dùng các gói đã khai báo trong main.tex; không thêm gói mới.
\chapter{NỀN TẢNG XỬ LÝ NGÔN NGỮ TIẾNG VIỆT}
\label{ch:ch2-nlp}

Chương này trình bày các khái niệm ngôn ngữ học và các quyết định tiền xử lý mà
hệ thống dựa vào. Mục tiêu không phải trình bày lý thuyết ngôn ngữ học đầy đủ, mà
là chỉ ra chính xác vì sao mỗi bước trong quy trình tiền xử lý tồn tại và nó ảnh
hưởng gì tới kết quả truy xuất.

\section{Âm tiết và từ trong tiếng Việt}

Tiếng Việt là ngôn ngữ đơn lập và có thanh điệu. Đơn vị ngữ âm cơ bản là
\emph{âm tiết}; trong chữ Quốc ngữ, các âm tiết được viết rời nhau bằng khoảng
trắng. Tuy nhiên đơn vị mang nghĩa dùng cho truy xuất lại là \emph{từ}, và một từ
thường gồm nhiều âm tiết. Ví dụ, ``trí tuệ nhân tạo'' là một cụm bốn âm tiết gắn
với một khái niệm, ``giáo dục'' là một từ hai âm tiết. Ranh giới từ không hiển
thị trực tiếp trên văn bản, nên cần một bộ tách từ để phục hồi ranh giới đó.

Sự lệch pha giữa cách viết (khoảng trắng) và đơn vị nghĩa (từ) là cái bẫy lớn
nhất khi lập chỉ mục tiếng Việt. Nếu coi mỗi âm tiết là một đơn vị chỉ mục, thì
``trí'', ``tuệ'', ``nhân'', ``tạo'' trở thành bốn chiều tách rời; khái niệm
``trí tuệ nhân tạo'' bị phân rã và tín hiệu phân biệt bị pha loãng bởi các âm tiết
rất phổ biến. Đây là lý do bước tách từ được đặt trước bước lập chỉ mục.

\section{Chuẩn hóa Unicode và dấu thanh}

Tiếng Việt dùng nhiều dấu thanh và dấu mũ, và cùng một chữ có thể được mã hóa theo
hai cách trong Unicode: dạng tổ hợp (một chữ cái cơ bản cộng dấu tổ hợp) và dạng
dựng sẵn (một điểm mã duy nhất). Hai chuỗi trông giống hệt nhau trên màn hình
nhưng khác nhau về dãy điểm mã sẽ không khớp nhau khi so sánh, khi tra từ điển
hay khi tính hàm băm. Vì vậy toàn bộ văn bản được chuẩn hóa về dạng dựng sẵn
Unicode NFC trước khi xử lý tiếp.

Bên cạnh chuẩn hóa tổ hợp dấu, quy trình còn hạ chữ hoa xuống chữ thường và loại
bỏ các thành phần không mang nội dung như thẻ đánh dấu, đường dẫn và địa chỉ thư
điện tử. Chi tiết các quy tắc này cùng lý do của từng quy tắc được trình bày ở
Chương~3, phần thiết kế tầng ngôn ngữ M1.

\section{Vì sao phải tách từ trước khi lập chỉ mục}

Biểu diễn túi từ (bag of words) bỏ qua thứ tự từ và coi mỗi đơn vị chỉ mục là một
chiều độc lập. Với ngôn ngữ dùng khoảng trắng phân tách từ, đơn vị đó là từ. Với
tiếng Việt, phải chủ động tạo ra đơn vị đó bằng tách từ. Tách từ đúng làm cho
``giáo dục'' là một chiều thay vì hai, và làm cho ``trí tuệ nhân tạo'' là một chiều
thay vì bốn. Cách làm này vừa giảm số chiều hiệu dụng, vừa tăng khả năng phân biệt
của các chiều còn lại.

Ngược lại, tách sai gây hại kép. Ghép nhầm hai đơn vị vốn độc lập làm mất một
chiều có nghĩa; cắt rời một từ ghép làm mất một chiều có nghĩa khác. Vì vậy tách
từ là bước cần được đo chất lượng bằng nhãn người chứ không chỉ bằng các ca kiểm
thử nhỏ.

\section{So khớp cực đại: FMM, BMM và BiMM}

Hướng tiếp cận được nhóm chọn là so khớp cực đại dựa trên từ điển, hoạt động hoàn
toàn ngoại tuyến và có vết thuật toán từng bước để kiểm chứng. Ý tưởng chung là
trượt một cửa sổ gồm tối đa $L$ âm tiết trên câu đã chuẩn hóa và chọn đoạn dài
nhất có mặt trong từ điển.

\begin{itemize}
  \item \textbf{FMM (so khớp tiến cực đại):} duyệt câu từ trái sang phải; tại mỗi
        vị trí, thử cửa sổ dài nhất rồi rút ngắn dần cho đến khi gặp một đoạn có
        trong từ điển.
  \item \textbf{BMM (so khớp lùi cực đại):} làm tương tự nhưng duyệt từ phải sang
        trái.
  \item \textbf{BiMM (so khớp hai chiều):} chạy cả hai hướng rồi chọn kết quả theo
        một bộ quy tắc phân xử.
\end{itemize}

Nhập nhằng là đặc trưng của phương pháp này. Với câu ``học sinh học sinh học'',
cả ``học sinh'' và ``sinh học'' đều là từ hợp lệ trong từ điển, nên FMM và BMM có
thể cho hai cách tách khác nhau trên cùng một câu. Bộ quy tắc phân xử của BiMM vì
vậy được định nghĩa theo thứ tự ưu tiên rõ ràng: chọn phương án có ít đơn vị hơn,
rồi ít âm tiết đơn hơn, rồi ít đơn vị ngoài từ điển hơn, và khi tất cả đều hòa thì
chọn BMM. Quy tắc này và cách kiểm chứng bằng từ điển fixture minh bạch được trình
bày ở Chương~3.

\section{Từ dừng, phủ định và bài toán lọc hai tầng}

Từ dừng là các từ chức năng xuất hiện rất phổ biến nhưng đóng góp ít cho việc phân
biệt tài liệu. Loại chúng đi giúp tăng trọng số tương đối của các từ mang nội dung
và giảm kích thước chỉ mục. Tuy nhiên có hai cạm bẫy.

Thứ nhất, một danh sách từ dừng cố định không tính đến miền dữ liệu. Từ rất phổ
biến trong kho tin tức thể thao có thể là từ hiếm và có giá trị trong một kho y
tế; ngược lại, có từ phổ biến chung nhưng lại tập trung ở vài chủ đề và mang
thông tin phân biệt. Vì vậy đồ án dùng hai tầng: tầng~1 loại hư từ ngữ pháp bằng
danh sách có nguồn và quy tắc biên soạn; tầng~2 dùng độ đo hình học TDV để phát
hiện từ dừng đặc thù của chính miền dữ liệu.

Thứ hai, phủ định mang thông tin điều kiện chân trị. Trong truy vấn ``bệnh không
truyền nhiễm'', ``không'' làm đảo nghĩa của cả cụm; loại ``không'' đi thì truy vấn
còn lại nói về ``bệnh truyền nhiễm'' -- nghĩa gần như ngược lại. Vì vậy theo sau
tầng lọc phải có một danh sách bảo vệ (guard) gỡ lại các từ phủ định và logic.
Đồ án giữ quy tắc: một từ được bảo vệ nếu chính nó nằm trong danh sách guard, hoặc
nếu một thành phần trong từ ghép của nó là guard được bảo vệ, nhờ đó các từ ghép
như ``không thể'' cũng được giữ nguyên vẹn.

Cần nói rõ giới hạn: giữ lại từ phủ định trong túi từ không biến mô hình thành bộ
suy luận phủ định theo ngữ nghĩa. Túi từ vẫn chỉ ghi nhận sự xuất hiện của từ
``không'', chứ không suy ra được quan hệ điều kiện chân trị phức tạp. Việc bảo vệ
guard chỉ ngăn chặn việc lọc làm mất dấu hiệu phủ định ở mức từ vựng.

\section{Từ điển tiếng Việt sử dụng}

Bước tách từ dựa trên một từ điển tiếng Việt đóng gói trong kho mã nguồn. Nhóm sử
dụng danh sách Viet74K \cite{viet74k}, một danh sách từ tiếng Việt mã nguồn mở,
được ghim ở một phiên bản (commit) cụ thể và đi kèm thông tin giấy phép. Từ điển
gồm 69.858 mục sau khi làm sạch; độ dài tối đa là 17 âm tiết, và khoảng 97,42\%
số mục có không quá bốn âm tiết. Con số độ dài tối đa dùng để đặt cửa sổ trượt
$L$ sao cho không bỏ sót mục nào, còn phân bố độ dài cho thấy phần lớn từ có độ
dài ngắn.

Đi kèm từ điển là một tệp thống kê mô tả số mục, phân bố độ dài và hàm băm của
tệp từ điển, để bất kỳ ai cũng có thể kiểm tra lại dữ liệu đã dùng. Chất lượng
tách từ phụ thuộc từ điển: đổi từ điển thì mọi số đo tách từ phải được đo lại, vì
cùng một câu có thể được tách khác đi.

\section{Từ vựng đến vector đặc trưng}

Sau tách từ và lọc từ dừng, mỗi tài liệu được biểu diễn thành một vector trọng số
trên từ vựng chung. Chương~3 trình bày công thức trọng số TF-IDF, cách chuẩn hóa
vector và cách tính tương đồng cosin. Trọng số không phải là số đếm thô: nó kết
hợp tần suất trong tài liệu với độ hiếm của từ trên toàn kho, nhờ đó các từ hiếm
nhưng phân biệt được đề cao hơn các từ phổ biến. Đây cũng là trực giác dẫn tới
tầng lọc thứ hai bằng TDV, bởi TDV đo trực tiếp tác động của việc bỏ một từ lên
khả năng phân biệt tài liệu, thay vì gián tiếp qua tần suất.
.
% Chỉ dùng các gói đã khai báo trong main.tex; không thêm gói mới.
\chapter{NỀN TẢNG XỬ LÝ NGÔN NGỮ TIẾNG VIỆT}
\label{ch:ch2-nlp}

Chương này trình bày các khái niệm ngôn ngữ học và các quyết định tiền xử lý mà
hệ thống dựa vào. Mục tiêu không phải trình bày lý thuyết ngôn ngữ học đầy đủ, mà
là chỉ ra chính xác vì sao mỗi bước trong quy trình tiền xử lý tồn tại và nó ảnh
hưởng gì tới kết quả truy xuất.

\section{Âm tiết và từ trong tiếng Việt}

Tiếng Việt là ngôn ngữ đơn lập và có thanh điệu. Đơn vị ngữ âm cơ bản là
\emph{âm tiết}; trong chữ Quốc ngữ, các âm tiết được viết rời nhau bằng khoảng
trắng. Tuy nhiên đơn vị mang nghĩa dùng cho truy xuất lại là \emph{từ}, và một từ
thường gồm nhiều âm tiết. Ví dụ, ``trí tuệ nhân tạo'' là một cụm bốn âm tiết gắn
với một khái niệm, ``giáo dục'' là một từ hai âm tiết. Ranh giới từ không hiển
thị trực tiếp trên văn bản, nên cần một bộ tách từ để phục hồi ranh giới đó.

Sự lệch pha giữa cách viết (khoảng trắng) và đơn vị nghĩa (từ) là cái bẫy lớn
nhất khi lập chỉ mục tiếng Việt. Nếu coi mỗi âm tiết là một đơn vị chỉ mục, thì
``trí'', ``tuệ'', ``nhân'', ``tạo'' trở thành bốn chiều tách rời; khái niệm
``trí tuệ nhân tạo'' bị phân rã và tín hiệu phân biệt bị pha loãng bởi các âm tiết
rất phổ biến. Đây là lý do bước tách từ được đặt trước bước lập chỉ mục.

\section{Chuẩn hóa Unicode và dấu thanh}

Tiếng Việt dùng nhiều dấu thanh và dấu mũ, và cùng một chữ có thể được mã hóa theo
hai cách trong Unicode: dạng tổ hợp (một chữ cái cơ bản cộng dấu tổ hợp) và dạng
dựng sẵn (một điểm mã duy nhất). Hai chuỗi trông giống hệt nhau trên màn hình
nhưng khác nhau về dãy điểm mã sẽ không khớp nhau khi so sánh, khi tra từ điển
hay khi tính hàm băm. Vì vậy toàn bộ văn bản được chuẩn hóa về dạng dựng sẵn
Unicode NFC trước khi xử lý tiếp.

Bên cạnh chuẩn hóa tổ hợp dấu, quy trình còn hạ chữ hoa xuống chữ thường và loại
bỏ các thành phần không mang nội dung như thẻ đánh dấu, đường dẫn và địa chỉ thư
điện tử. Chi tiết các quy tắc này cùng lý do của từng quy tắc được trình bày ở
Chương~3, phần thiết kế tầng ngôn ngữ M1.

\section{Vì sao phải tách từ trước khi lập chỉ mục}

Biểu diễn túi từ (bag of words) bỏ qua thứ tự từ và coi mỗi đơn vị chỉ mục là một
chiều độc lập. Với ngôn ngữ dùng khoảng trắng phân tách từ, đơn vị đó là từ. Với
tiếng Việt, phải chủ động tạo ra đơn vị đó bằng tách từ. Tách từ đúng làm cho
``giáo dục'' là một chiều thay vì hai, và làm cho ``trí tuệ nhân tạo'' là một chiều
thay vì bốn. Cách làm này vừa giảm số chiều hiệu dụng, vừa tăng khả năng phân biệt
của các chiều còn lại.

Ngược lại, tách sai gây hại kép. Ghép nhầm hai đơn vị vốn độc lập làm mất một
chiều có nghĩa; cắt rời một từ ghép làm mất một chiều có nghĩa khác. Vì vậy tách
từ là bước cần được đo chất lượng bằng nhãn người chứ không chỉ bằng các ca kiểm
thử nhỏ.

\section{So khớp cực đại: FMM, BMM và BiMM}

Hướng tiếp cận được nhóm chọn là so khớp cực đại dựa trên từ điển, hoạt động hoàn
toàn ngoại tuyến và có vết thuật toán từng bước để kiểm chứng. Ý tưởng chung là
trượt một cửa sổ gồm tối đa $L$ âm tiết trên câu đã chuẩn hóa và chọn đoạn dài
nhất có mặt trong từ điển.

\begin{itemize}
  \item \textbf{FMM (so khớp tiến cực đại):} duyệt câu từ trái sang phải; tại mỗi
        vị trí, thử cửa sổ dài nhất rồi rút ngắn dần cho đến khi gặp một đoạn có
        trong từ điển.
  \item \textbf{BMM (so khớp lùi cực đại):} làm tương tự nhưng duyệt từ phải sang
        trái.
  \item \textbf{BiMM (so khớp hai chiều):} chạy cả hai hướng rồi chọn kết quả theo
        một bộ quy tắc phân xử.
\end{itemize}

Nhập nhằng là đặc trưng của phương pháp này. Với câu ``học sinh học sinh học'',
cả ``học sinh'' và ``sinh học'' đều là từ hợp lệ trong từ điển, nên FMM và BMM có
thể cho hai cách tách khác nhau trên cùng một câu. Bộ quy tắc phân xử của BiMM vì
vậy được định nghĩa theo thứ tự ưu tiên rõ ràng: chọn phương án có ít đơn vị hơn,
rồi ít âm tiết đơn hơn, rồi ít đơn vị ngoài từ điển hơn, và khi tất cả đều hòa thì
chọn BMM. Quy tắc này và cách kiểm chứng bằng từ điển fixture minh bạch được trình
bày ở Chương~3.

\section{Từ dừng, phủ định và bài toán lọc hai tầng}

Từ dừng là các từ chức năng xuất hiện rất phổ biến nhưng đóng góp ít cho việc phân
biệt tài liệu. Loại chúng đi giúp tăng trọng số tương đối của các từ mang nội dung
và giảm kích thước chỉ mục. Tuy nhiên có hai cạm bẫy.

Thứ nhất, một danh sách từ dừng cố định không tính đến miền dữ liệu. Từ rất phổ
biến trong kho tin tức thể thao có thể là từ hiếm và có giá trị trong một kho y
tế; ngược lại, có từ phổ biến chung nhưng lại tập trung ở vài chủ đề và mang
thông tin phân biệt. Vì vậy đồ án dùng hai tầng: tầng~1 loại hư từ ngữ pháp bằng
danh sách có nguồn và quy tắc biên soạn; tầng~2 dùng độ đo hình học TDV để phát
hiện từ dừng đặc thù của chính miền dữ liệu.

Thứ hai, phủ định mang thông tin điều kiện chân trị. Trong truy vấn ``bệnh không
truyền nhiễm'', ``không'' làm đảo nghĩa của cả cụm; loại ``không'' đi thì truy vấn
còn lại nói về ``bệnh truyền nhiễm'' -- nghĩa gần như ngược lại. Vì vậy theo sau
tầng lọc phải có một danh sách bảo vệ (guard) gỡ lại các từ phủ định và logic.
Đồ án giữ quy tắc: một từ được bảo vệ nếu chính nó nằm trong danh sách guard, hoặc
nếu một thành phần trong từ ghép của nó là guard được bảo vệ, nhờ đó các từ ghép
như ``không thể'' cũng được giữ nguyên vẹn.

Cần nói rõ giới hạn: giữ lại từ phủ định trong túi từ không biến mô hình thành bộ
suy luận phủ định theo ngữ nghĩa. Túi từ vẫn chỉ ghi nhận sự xuất hiện của từ
``không'', chứ không suy ra được quan hệ điều kiện chân trị phức tạp. Việc bảo vệ
guard chỉ ngăn chặn việc lọc làm mất dấu hiệu phủ định ở mức từ vựng.

\section{Từ điển tiếng Việt sử dụng}

Bước tách từ dựa trên một từ điển tiếng Việt đóng gói trong kho mã nguồn. Nhóm sử
dụng danh sách Viet74K \cite{viet74k}, một danh sách từ tiếng Việt mã nguồn mở,
được ghim ở một phiên bản (commit) cụ thể và đi kèm thông tin giấy phép. Từ điển
gồm 69.858 mục sau khi làm sạch; độ dài tối đa là 17 âm tiết, và khoảng 97,42\%
số mục có không quá bốn âm tiết. Con số độ dài tối đa dùng để đặt cửa sổ trượt
$L$ sao cho không bỏ sót mục nào, còn phân bố độ dài cho thấy phần lớn từ có độ
dài ngắn.

Đi kèm từ điển là một tệp thống kê mô tả số mục, phân bố độ dài và hàm băm của
tệp từ điển, để bất kỳ ai cũng có thể kiểm tra lại dữ liệu đã dùng. Chất lượng
tách từ phụ thuộc từ điển: đổi từ điển thì mọi số đo tách từ phải được đo lại, vì
cùng một câu có thể được tách khác đi.

\section{Từ vựng đến vector đặc trưng}

Sau tách từ và lọc từ dừng, mỗi tài liệu được biểu diễn thành một vector trọng số
trên từ vựng chung. Chương~3 trình bày công thức trọng số TF-IDF, cách chuẩn hóa
vector và cách tính tương đồng cosin. Trọng số không phải là số đếm thô: nó kết
hợp tần suất trong tài liệu với độ hiếm của từ trên toàn kho, nhờ đó các từ hiếm
nhưng phân biệt được đề cao hơn các từ phổ biến. Đây cũng là trực giác dẫn tới
tầng lọc thứ hai bằng TDV, bởi TDV đo trực tiếp tác động của việc bỏ một từ lên
khả năng phân biệt tài liệu, thay vì gián tiếp qua tần suất.
.
% Chỉ dùng các gói đã khai báo trong main.tex; không thêm gói mới.
\chapter{NỀN TẢNG XỬ LÝ NGÔN NGỮ TIẾNG VIỆT}
\label{ch:ch2-nlp}

Chương này trình bày các khái niệm ngôn ngữ học và các quyết định tiền xử lý mà
hệ thống dựa vào. Mục tiêu không phải trình bày lý thuyết ngôn ngữ học đầy đủ, mà
là chỉ ra chính xác vì sao mỗi bước trong quy trình tiền xử lý tồn tại và nó ảnh
hưởng gì tới kết quả truy xuất.

\section{Âm tiết và từ trong tiếng Việt}

Tiếng Việt là ngôn ngữ đơn lập và có thanh điệu. Đơn vị ngữ âm cơ bản là
\emph{âm tiết}; trong chữ Quốc ngữ, các âm tiết được viết rời nhau bằng khoảng
trắng. Tuy nhiên đơn vị mang nghĩa dùng cho truy xuất lại là \emph{từ}, và một từ
thường gồm nhiều âm tiết. Ví dụ, ``trí tuệ nhân tạo'' là một cụm bốn âm tiết gắn
với một khái niệm, ``giáo dục'' là một từ hai âm tiết. Ranh giới từ không hiển
thị trực tiếp trên văn bản, nên cần một bộ tách từ để phục hồi ranh giới đó.

Sự lệch pha giữa cách viết (khoảng trắng) và đơn vị nghĩa (từ) là cái bẫy lớn
nhất khi lập chỉ mục tiếng Việt. Nếu coi mỗi âm tiết là một đơn vị chỉ mục, thì
``trí'', ``tuệ'', ``nhân'', ``tạo'' trở thành bốn chiều tách rời; khái niệm
``trí tuệ nhân tạo'' bị phân rã và tín hiệu phân biệt bị pha loãng bởi các âm tiết
rất phổ biến. Đây là lý do bước tách từ được đặt trước bước lập chỉ mục.

\section{Chuẩn hóa Unicode và dấu thanh}

Tiếng Việt dùng nhiều dấu thanh và dấu mũ, và cùng một chữ có thể được mã hóa theo
hai cách trong Unicode: dạng tổ hợp (một chữ cái cơ bản cộng dấu tổ hợp) và dạng
dựng sẵn (một điểm mã duy nhất). Hai chuỗi trông giống hệt nhau trên màn hình
nhưng khác nhau về dãy điểm mã sẽ không khớp nhau khi so sánh, khi tra từ điển
hay khi tính hàm băm. Vì vậy toàn bộ văn bản được chuẩn hóa về dạng dựng sẵn
Unicode NFC trước khi xử lý tiếp.

Bên cạnh chuẩn hóa tổ hợp dấu, quy trình còn hạ chữ hoa xuống chữ thường và loại
bỏ các thành phần không mang nội dung như thẻ đánh dấu, đường dẫn và địa chỉ thư
điện tử. Chi tiết các quy tắc này cùng lý do của từng quy tắc được trình bày ở
Chương~3, phần thiết kế tầng ngôn ngữ M1.

\section{Vì sao phải tách từ trước khi lập chỉ mục}

Biểu diễn túi từ (bag of words) bỏ qua thứ tự từ và coi mỗi đơn vị chỉ mục là một
chiều độc lập. Với ngôn ngữ dùng khoảng trắng phân tách từ, đơn vị đó là từ. Với
tiếng Việt, phải chủ động tạo ra đơn vị đó bằng tách từ. Tách từ đúng làm cho
``giáo dục'' là một chiều thay vì hai, và làm cho ``trí tuệ nhân tạo'' là một chiều
thay vì bốn. Cách làm này vừa giảm số chiều hiệu dụng, vừa tăng khả năng phân biệt
của các chiều còn lại.

Ngược lại, tách sai gây hại kép. Ghép nhầm hai đơn vị vốn độc lập làm mất một
chiều có nghĩa; cắt rời một từ ghép làm mất một chiều có nghĩa khác. Vì vậy tách
từ là bước cần được đo chất lượng bằng nhãn người chứ không chỉ bằng các ca kiểm
thử nhỏ.

\section{So khớp cực đại: FMM, BMM và BiMM}

Hướng tiếp cận được nhóm chọn là so khớp cực đại dựa trên từ điển, hoạt động hoàn
toàn ngoại tuyến và có vết thuật toán từng bước để kiểm chứng. Ý tưởng chung là
trượt một cửa sổ gồm tối đa $L$ âm tiết trên câu đã chuẩn hóa và chọn đoạn dài
nhất có mặt trong từ điển.

\begin{itemize}
  \item \textbf{FMM (so khớp tiến cực đại):} duyệt câu từ trái sang phải; tại mỗi
        vị trí, thử cửa sổ dài nhất rồi rút ngắn dần cho đến khi gặp một đoạn có
        trong từ điển.
  \item \textbf{BMM (so khớp lùi cực đại):} làm tương tự nhưng duyệt từ phải sang
        trái.
  \item \textbf{BiMM (so khớp hai chiều):} chạy cả hai hướng rồi chọn kết quả theo
        một bộ quy tắc phân xử.
\end{itemize}

Nhập nhằng là đặc trưng của phương pháp này. Với câu ``học sinh học sinh học'',
cả ``học sinh'' và ``sinh học'' đều là từ hợp lệ trong từ điển, nên FMM và BMM có
thể cho hai cách tách khác nhau trên cùng một câu. Bộ quy tắc phân xử của BiMM vì
vậy được định nghĩa theo thứ tự ưu tiên rõ ràng: chọn phương án có ít đơn vị hơn,
rồi ít âm tiết đơn hơn, rồi ít đơn vị ngoài từ điển hơn, và khi tất cả đều hòa thì
chọn BMM. Quy tắc này và cách kiểm chứng bằng từ điển fixture minh bạch được trình
bày ở Chương~3.

\section{Từ dừng, phủ định và bài toán lọc hai tầng}

Từ dừng là các từ chức năng xuất hiện rất phổ biến nhưng đóng góp ít cho việc phân
biệt tài liệu. Loại chúng đi giúp tăng trọng số tương đối của các từ mang nội dung
và giảm kích thước chỉ mục. Tuy nhiên có hai cạm bẫy.

Thứ nhất, một danh sách từ dừng cố định không tính đến miền dữ liệu. Từ rất phổ
biến trong kho tin tức thể thao có thể là từ hiếm và có giá trị trong một kho y
tế; ngược lại, có từ phổ biến chung nhưng lại tập trung ở vài chủ đề và mang
thông tin phân biệt. Vì vậy đồ án dùng hai tầng: tầng~1 loại hư từ ngữ pháp bằng
danh sách có nguồn và quy tắc biên soạn; tầng~2 dùng độ đo hình học TDV để phát
hiện từ dừng đặc thù của chính miền dữ liệu.

Thứ hai, phủ định mang thông tin điều kiện chân trị. Trong truy vấn ``bệnh không
truyền nhiễm'', ``không'' làm đảo nghĩa của cả cụm; loại ``không'' đi thì truy vấn
còn lại nói về ``bệnh truyền nhiễm'' -- nghĩa gần như ngược lại. Vì vậy theo sau
tầng lọc phải có một danh sách bảo vệ (guard) gỡ lại các từ phủ định và logic.
Đồ án giữ quy tắc: một từ được bảo vệ nếu chính nó nằm trong danh sách guard, hoặc
nếu một thành phần trong từ ghép của nó là guard được bảo vệ, nhờ đó các từ ghép
như ``không thể'' cũng được giữ nguyên vẹn.

Cần nói rõ giới hạn: giữ lại từ phủ định trong túi từ không biến mô hình thành bộ
suy luận phủ định theo ngữ nghĩa. Túi từ vẫn chỉ ghi nhận sự xuất hiện của từ
``không'', chứ không suy ra được quan hệ điều kiện chân trị phức tạp. Việc bảo vệ
guard chỉ ngăn chặn việc lọc làm mất dấu hiệu phủ định ở mức từ vựng.

\section{Từ điển tiếng Việt sử dụng}

Bước tách từ dựa trên một từ điển tiếng Việt đóng gói trong kho mã nguồn. Nhóm sử
dụng danh sách Viet74K \cite{viet74k}, một danh sách từ tiếng Việt mã nguồn mở,
được ghim ở một phiên bản (commit) cụ thể và đi kèm thông tin giấy phép. Từ điển
gồm 69.858 mục sau khi làm sạch; độ dài tối đa là 17 âm tiết, và khoảng 97,42\%
số mục có không quá bốn âm tiết. Con số độ dài tối đa dùng để đặt cửa sổ trượt
$L$ sao cho không bỏ sót mục nào, còn phân bố độ dài cho thấy phần lớn từ có độ
dài ngắn.

Đi kèm từ điển là một tệp thống kê mô tả số mục, phân bố độ dài và hàm băm của
tệp từ điển, để bất kỳ ai cũng có thể kiểm tra lại dữ liệu đã dùng. Chất lượng
tách từ phụ thuộc từ điển: đổi từ điển thì mọi số đo tách từ phải được đo lại, vì
cùng một câu có thể được tách khác đi.

\section{Từ vựng đến vector đặc trưng}

Sau tách từ và lọc từ dừng, mỗi tài liệu được biểu diễn thành một vector trọng số
trên từ vựng chung. Chương~3 trình bày công thức trọng số TF-IDF, cách chuẩn hóa
vector và cách tính tương đồng cosin. Trọng số không phải là số đếm thô: nó kết
hợp tần suất trong tài liệu với độ hiếm của từ trên toàn kho, nhờ đó các từ hiếm
nhưng phân biệt được đề cao hơn các từ phổ biến. Đây cũng là trực giác dẫn tới
tầng lọc thứ hai bằng TDV, bởi TDV đo trực tiếp tác động của việc bỏ một từ lên
khả năng phân biệt tài liệu, thay vì gián tiếp qua tần suất.
.
% Chỉ dùng các gói đã khai báo trong main.tex; không thêm gói mới.
\chapter{NỀN TẢNG XỬ LÝ NGÔN NGỮ TIẾNG VIỆT}
\label{ch:ch2-nlp}

Chương này trình bày các khái niệm ngôn ngữ học và các quyết định tiền xử lý mà
hệ thống dựa vào. Mục tiêu không phải trình bày lý thuyết ngôn ngữ học đầy đủ, mà
là chỉ ra chính xác vì sao mỗi bước trong quy trình tiền xử lý tồn tại và nó ảnh
hưởng gì tới kết quả truy xuất.

\section{Âm tiết và từ trong tiếng Việt}

Tiếng Việt là ngôn ngữ đơn lập và có thanh điệu. Đơn vị ngữ âm cơ bản là
\emph{âm tiết}; trong chữ Quốc ngữ, các âm tiết được viết rời nhau bằng khoảng
trắng. Tuy nhiên đơn vị mang nghĩa dùng cho truy xuất lại là \emph{từ}, và một từ
thường gồm nhiều âm tiết. Ví dụ, ``trí tuệ nhân tạo'' là một cụm bốn âm tiết gắn
với một khái niệm, ``giáo dục'' là một từ hai âm tiết. Ranh giới từ không hiển
thị trực tiếp trên văn bản, nên cần một bộ tách từ để phục hồi ranh giới đó.

Sự lệch pha giữa cách viết (khoảng trắng) và đơn vị nghĩa (từ) là cái bẫy lớn
nhất khi lập chỉ mục tiếng Việt. Nếu coi mỗi âm tiết là một đơn vị chỉ mục, thì
``trí'', ``tuệ'', ``nhân'', ``tạo'' trở thành bốn chiều tách rời; khái niệm
``trí tuệ nhân tạo'' bị phân rã và tín hiệu phân biệt bị pha loãng bởi các âm tiết
rất phổ biến. Đây là lý do bước tách từ được đặt trước bước lập chỉ mục.

\section{Chuẩn hóa Unicode và dấu thanh}

Tiếng Việt dùng nhiều dấu thanh và dấu mũ, và cùng một chữ có thể được mã hóa theo
hai cách trong Unicode: dạng tổ hợp (một chữ cái cơ bản cộng dấu tổ hợp) và dạng
dựng sẵn (một điểm mã duy nhất). Hai chuỗi trông giống hệt nhau trên màn hình
nhưng khác nhau về dãy điểm mã sẽ không khớp nhau khi so sánh, khi tra từ điển
hay khi tính hàm băm. Vì vậy toàn bộ văn bản được chuẩn hóa về dạng dựng sẵn
Unicode NFC trước khi xử lý tiếp.

Bên cạnh chuẩn hóa tổ hợp dấu, quy trình còn hạ chữ hoa xuống chữ thường và loại
bỏ các thành phần không mang nội dung như thẻ đánh dấu, đường dẫn và địa chỉ thư
điện tử. Chi tiết các quy tắc này cùng lý do của từng quy tắc được trình bày ở
Chương~3, phần thiết kế tầng ngôn ngữ M1.

\section{Vì sao phải tách từ trước khi lập chỉ mục}

Biểu diễn túi từ (bag of words) bỏ qua thứ tự từ và coi mỗi đơn vị chỉ mục là một
chiều độc lập. Với ngôn ngữ dùng khoảng trắng phân tách từ, đơn vị đó là từ. Với
tiếng Việt, phải chủ động tạo ra đơn vị đó bằng tách từ. Tách từ đúng làm cho
``giáo dục'' là một chiều thay vì hai, và làm cho ``trí tuệ nhân tạo'' là một chiều
thay vì bốn. Cách làm này vừa giảm số chiều hiệu dụng, vừa tăng khả năng phân biệt
của các chiều còn lại.

Ngược lại, tách sai gây hại kép. Ghép nhầm hai đơn vị vốn độc lập làm mất một
chiều có nghĩa; cắt rời một từ ghép làm mất một chiều có nghĩa khác. Vì vậy tách
từ là bước cần được đo chất lượng bằng nhãn người chứ không chỉ bằng các ca kiểm
thử nhỏ.

\section{So khớp cực đại: FMM, BMM và BiMM}

Hướng tiếp cận được nhóm chọn là so khớp cực đại dựa trên từ điển, hoạt động hoàn
toàn ngoại tuyến và có vết thuật toán từng bước để kiểm chứng. Ý tưởng chung là
trượt một cửa sổ gồm tối đa $L$ âm tiết trên câu đã chuẩn hóa và chọn đoạn dài
nhất có mặt trong từ điển.

\begin{itemize}
  \item \textbf{FMM (so khớp tiến cực đại):} duyệt câu từ trái sang phải; tại mỗi
        vị trí, thử cửa sổ dài nhất rồi rút ngắn dần cho đến khi gặp một đoạn có
        trong từ điển.
  \item \textbf{BMM (so khớp lùi cực đại):} làm tương tự nhưng duyệt từ phải sang
        trái.
  \item \textbf{BiMM (so khớp hai chiều):} chạy cả hai hướng rồi chọn kết quả theo
        một bộ quy tắc phân xử.
\end{itemize}

Nhập nhằng là đặc trưng của phương pháp này. Với câu ``học sinh học sinh học'',
cả ``học sinh'' và ``sinh học'' đều là từ hợp lệ trong từ điển, nên FMM và BMM có
thể cho hai cách tách khác nhau trên cùng một câu. Bộ quy tắc phân xử của BiMM vì
vậy được định nghĩa theo thứ tự ưu tiên rõ ràng: chọn phương án có ít đơn vị hơn,
rồi ít âm tiết đơn hơn, rồi ít đơn vị ngoài từ điển hơn, và khi tất cả đều hòa thì
chọn BMM. Quy tắc này và cách kiểm chứng bằng từ điển fixture minh bạch được trình
bày ở Chương~3.

\section{Từ dừng, phủ định và bài toán lọc hai tầng}

Từ dừng là các từ chức năng xuất hiện rất phổ biến nhưng đóng góp ít cho việc phân
biệt tài liệu. Loại chúng đi giúp tăng trọng số tương đối của các từ mang nội dung
và giảm kích thước chỉ mục. Tuy nhiên có hai cạm bẫy.

Thứ nhất, một danh sách từ dừng cố định không tính đến miền dữ liệu. Từ rất phổ
biến trong kho tin tức thể thao có thể là từ hiếm và có giá trị trong một kho y
tế; ngược lại, có từ phổ biến chung nhưng lại tập trung ở vài chủ đề và mang
thông tin phân biệt. Vì vậy đồ án dùng hai tầng: tầng~1 loại hư từ ngữ pháp bằng
danh sách có nguồn và quy tắc biên soạn; tầng~2 dùng độ đo hình học TDV để phát
hiện từ dừng đặc thù của chính miền dữ liệu.

Thứ hai, phủ định mang thông tin điều kiện chân trị. Trong truy vấn ``bệnh không
truyền nhiễm'', ``không'' làm đảo nghĩa của cả cụm; loại ``không'' đi thì truy vấn
còn lại nói về ``bệnh truyền nhiễm'' -- nghĩa gần như ngược lại. Vì vậy theo sau
tầng lọc phải có một danh sách bảo vệ (guard) gỡ lại các từ phủ định và logic.
Đồ án giữ quy tắc: một từ được bảo vệ nếu chính nó nằm trong danh sách guard, hoặc
nếu một thành phần trong từ ghép của nó là guard được bảo vệ, nhờ đó các từ ghép
như ``không thể'' cũng được giữ nguyên vẹn.

Cần nói rõ giới hạn: giữ lại từ phủ định trong túi từ không biến mô hình thành bộ
suy luận phủ định theo ngữ nghĩa. Túi từ vẫn chỉ ghi nhận sự xuất hiện của từ
``không'', chứ không suy ra được quan hệ điều kiện chân trị phức tạp. Việc bảo vệ
guard chỉ ngăn chặn việc lọc làm mất dấu hiệu phủ định ở mức từ vựng.

\section{Từ điển tiếng Việt sử dụng}

Bước tách từ dựa trên một từ điển tiếng Việt đóng gói trong kho mã nguồn. Nhóm sử
dụng danh sách Viet74K \cite{viet74k}, một danh sách từ tiếng Việt mã nguồn mở,
được ghim ở một phiên bản (commit) cụ thể và đi kèm thông tin giấy phép. Từ điển
gồm 69.858 mục sau khi làm sạch; độ dài tối đa là 17 âm tiết, và khoảng 97,42\%
số mục có không quá bốn âm tiết. Con số độ dài tối đa dùng để đặt cửa sổ trượt
$L$ sao cho không bỏ sót mục nào, còn phân bố độ dài cho thấy phần lớn từ có độ
dài ngắn.

Đi kèm từ điển là một tệp thống kê mô tả số mục, phân bố độ dài và hàm băm của
tệp từ điển, để bất kỳ ai cũng có thể kiểm tra lại dữ liệu đã dùng. Chất lượng
tách từ phụ thuộc từ điển: đổi từ điển thì mọi số đo tách từ phải được đo lại, vì
cùng một câu có thể được tách khác đi.

\section{Từ vựng đến vector đặc trưng}

Sau tách từ và lọc từ dừng, mỗi tài liệu được biểu diễn thành một vector trọng số
trên từ vựng chung. Chương~3 trình bày công thức trọng số TF-IDF, cách chuẩn hóa
vector và cách tính tương đồng cosin. Trọng số không phải là số đếm thô: nó kết
hợp tần suất trong tài liệu với độ hiếm của từ trên toàn kho, nhờ đó các từ hiếm
nhưng phân biệt được đề cao hơn các từ phổ biến. Đây cũng là trực giác dẫn tới
tầng lọc thứ hai bằng TDV, bởi TDV đo trực tiếp tác động của việc bỏ một từ lên
khả năng phân biệt tài liệu, thay vì gián tiếp qua tần suất.
